\documentclass[10pt]{article}

\usepackage[letterpaper,margin=0.85in]{geometry}
\usepackage{amsmath}
\usepackage{array}
\usepackage{booktabs}
\usepackage{enumitem}
\usepackage{hyperref}

\hypersetup{
  colorlinks=true,
  linkcolor=black,
  citecolor=black,
  urlcolor=blue
}

\setlength{\parindent}{1.2em}
\setlength{\parskip}{0.25em}
\emergencystretch=2em
\setlist[itemize]{leftmargin=1.5em}
\setlist[enumerate]{leftmargin=1.5em}

\title{\textbf{Physics-Backed Dynamic Policy Search for Pilot-Production Viability}\\
\large AIAA-Style Technical Note}
\author{Tyler K. B.\\
\normalsize Branch: \texttt{codex-viability-prototype}}
\date{June 2026}

\begin{document}

\maketitle

\begin{abstract}
This note documents the viability work on the
\texttt{codex-viability-prototype} branch. The workflow now uses
physics-backed direct evaluation as the authority for feasibility claims and
casts dynamic policy design as a finite-horizon nonlinear optimal control
problem. The implemented policy class uses three open-loop epochs. Each epoch
sets seven controls: annual intake, retention, UTE, PAA, maximum manning, FLUG
quota, and IPUG quota. A Gaussian-process model is used only to propose
candidates; every reported result is checked with direct physics. The closeout
campaign evaluated 2158 unique three- and five-epoch direct-physics schedules
over a 60-phase horizon. No feasible policy was found. The best policy had
$\phi=5.83075$ and required simultaneous relaxation of the total-pilot window,
WG RAP, and FL RAP constraints. A local finite-difference diagnostic shows that
the binding constraints respond to different controls, and in some cases move in
opposing directions. The result does not prove global infeasibility, but it gives
a concrete nearest miss and a defensible basis for requirements trade studies.
\end{abstract}

\section*{Nomenclature}

\begin{tabular}{>{\raggedright\arraybackslash}p{0.18\linewidth}p{0.76\linewidth}}
$x_k$ & Compressed force and training state at phase $k$, represented by simulator history: total and line pilots, qualification mix, staff counts, upgrade and carryover burden, and RAP shortfalls. \\
$u_k$ & Policy control vector applied at phase or epoch $k$: annual intake, retention, UTE, PAA, max manning, FLUG quota, and IPUG quota. \\
$f_k$ & Physics-backed simulator transition map from $x_k$ and $u_k$ to $x_{k+1}$. \\
$g_i(x)$ & Constraint function for target pilot inventory, RAP shortfalls, staff counts, or experience requirements. A feasible policy satisfies $g_i(x) \le 0$ for all enabled constraints. \\
$\phi$ & Maximum normalized constraint violation; $\phi \le 0$ denotes direct-verified feasibility. \\
RAP & Ready aircrew program requirement proxy used by the viability constraints. \\
WG, FL, IP & Wingman, flight lead, and instructor pilot qualification categories. \\
DOE & Design of experiments. \\
GPR & Gaussian-process regression. \\
\end{tabular}

\section{Introduction}

The earlier viability workflow used constant policies, static DOE studies,
surrogate screening, and two-dimensional feasible-envelope plots. That was a
necessary first analysis pass, but it was not enough for policy design. A
constant policy cannot represent a phased recovery strategy. A surrogate-only
screen cannot support final feasibility claims. Static envelope plots also
become difficult to use once several policy levers and constraints are active.

This branch addresses those issues directly. The long-horizon evaluator can now
use the single-phase physics allocator instead of the pretrained internal MLP.
The direct-physics path has been benchmarked and optimized enough to support
parallel candidate batches. A new structured dynamic-policy workflow then
searches over three open-loop epochs while keeping direct physics as the source
of truth.

\section{Physics-Backed Viability Baseline}

The evaluator now separates two levels of ``direct'' evaluation. A direct
long-horizon viability run bypasses the outer feasibility surrogate and runs
\texttt{CAFSimulation} for a selected policy. A physics-backed long-horizon run
also bypasses the internal MLP phase surrogate and calls the single-phase
physics allocator inside each simulated phase. The runs reported here use
\texttt{model.phase\_backend: physics}.

\subsection{Why the Internal MLP Is Not Authoritative}

The pretrained internal phase MLP was audited before this dynamic-policy work
was started. On the configured 16-output target set, a Sobol holdout audit gave
an all-target $R^2$ of about 0.595. The sortie-rate group, which is central to
the RAP constraints, was worse: its $R^2$ was about 0.244. That level of error
is acceptable for historical context and debugging, but it is not acceptable
for final policy claims.

The follow-on surrogate work improved the screening model substantially. A
4096-row shared ARD Matern $\nu=2.5$ GPR reached holdout $R^2$ values of about
0.933 over all targets, 0.881 over sortie rates, and 0.949 over simulator-rate
targets. Those results made the GPR useful for candidate generation. They did
not change the decision rule: final feasibility claims in this branch require
direct physics.

The performance work changed the search strategy. Early direct
physics runs were slow enough to make broad search prohibitively expensive. After
allocator and event-capacity optimizations, a documented one-year, eight-policy
probe improved from 23.8 s with one worker and 5.8 s with eight workers to
4.1 s with one worker and 1.8 s with eight workers. That made direct physics
fast enough for DOE, candidate verification, and moderate structured search. The
baseline used here is therefore:
\begin{quote}
Use physics-backed direct evaluation with \texttt{--workers} for final
viability claims. Use surrogates only for guidance, screening, and candidate
proposal.
\end{quote}

\section{Finite-Horizon Dynamic-Policy Formulation}

Dynamic policy design is treated as an open-loop finite-horizon nonlinear
optimal control problem. The horizon contains 60 simulation phases. The
implemented policy class uses three equal epochs instead of a fully free
per-phase schedule. Each epoch holds seven controls constant over its phase
block. The resulting decision vector has 21 dimensions:
\begin{equation}
U =
\left[
u_1^\top, u_2^\top, u_3^\top
\right]^\top,
\quad
u_e =
\begin{bmatrix}
\text{intake}_e &
\text{retention}_e &
\text{UTE}_e &
\text{PAA}_e &
\text{manning}_e &
\text{FLUG}_e &
\text{IPUG}_e
\end{bmatrix}^\top .
\end{equation}

The simulator defines the nonlinear transition map,
\begin{equation}
x_{k+1} = f_k(x_k,u_{e(k)}),
\end{equation}
where $e(k)$ maps phase $k$ to one of the three policy epochs. The objective is
to minimize the maximum normalized constraint violation,
\begin{equation}
\min_U \; \phi(U)
       = \min_U \; \max_i \frac{g_i(x_{0:N};U)}{s_i},
\end{equation}
where $s_i$ is the configured constraint scale. A policy is not labeled feasible
unless direct physics returns $\phi \le 0$.

\section{Implementation Summary}

The implementation keeps the dynamic search inside the \texttt{src/viability}
boundary. The main changes are:
\begin{itemize}
  \item \texttt{src/viability/dynamic\_policy.py}: defines epoch-policy
  schedules, phase-to-epoch mapping, flattened policy feature names, and
  unit-cube conversion for bounded policy variables.
  \item \texttt{src/viability/evaluator.py}: adds direct evaluation of
  \texttt{EpochPolicySchedule} objects by applying the active epoch policy
  before each simulator phase.
  \item \texttt{src/viability/dynamic\_search.py}: adds Sobol initialization,
  interpolated structured heuristic seeds, GPR fitting, surrogate-guided
  candidate proposal, direct physics verification, targeted refinement, and
  local finite-difference diagnostics.
  \item \texttt{src/viability/relaxation.py}: adds observed-search relaxation
  summaries, Pareto-front extraction, and constraint-set relaxation tables.
  \item \texttt{src/viability/cli.py}: adds \texttt{dynamic-search} and
  \texttt{dynamic-diagnose} commands, plus \texttt{dynamic-refine} and
  \texttt{dynamic-relaxation-study} for closeout analysis.
  \item \texttt{viability\_dashboard.py}: defaults to dynamic search results,
  exposes relaxation artifacts, and keeps legacy constant-policy sliders as a
  secondary mode.
  \item \texttt{tests/test\_viability\_dynamic\_policy.py} and
  \texttt{tests/test\_viability\_dynamic\_search.py}: cover epoch mapping,
  integer rounding, dynamic simulation policy application, artifact contracts,
  and diagnostic report generation.
\end{itemize}

Generated run artifacts stay under ignored \path{outputs/viability/...}; no
generated result files are tracked.

\section{Search Procedure}

The search used the physics-backed configuration in
\path{outputs/viability/dynamic_policy_search/config_physics.yaml}. The
initial three-epoch procedure was:
\begin{enumerate}
  \item Generate three-epoch schedules from structured heuristic seeds and a
  Sobol initial design.
  \item Evaluate all initial schedules with direct physics using eight workers.
  \item Fit a Matern-$\nu=1.5$ GPR surrogate to direct-physics $\phi$ values in
  normalized control space.
  \item Score a large Sobol candidate pool with a lower-confidence-bound rule
  and add one differential-evolution optimum from the surrogate.
  \item Direct-verify the top optimizer-proposed schedules with the physics
  backend.
  \item Run finite-difference perturbations around the best direct-verified
  schedule.
\end{enumerate}

That run evaluated 519 initial schedules and 96 optimizer-proposed schedules,
for 615 direct physics evaluations. Closeout then added two analysis passes:
\begin{itemize}
  \item a three-epoch targeted refinement around the nearest misses, with 2048
  local Sobol perturbations, a 131072-point GPR-screened pool, and 256 new
  direct checks;
  \item a five-epoch open-loop search with interpolated named heuristic seeds,
  1024 Sobol initial samples, a 131072-point GPR-screened pool, and 256
  optimizer checks.
\end{itemize}

The five-epoch search did not improve the best $\phi$ or change the binding
tradeoff story, so the conditional 10-epoch search was not run in this branch.

\section{Results}

No feasible policy was found in the tested structured open-loop policy classes.
Table
\ref{tab:search-summary} gives the run summary.

\begin{table}[htbp]
\centering
\caption{Dynamic-policy closeout search summary.}
\label{tab:search-summary}
\begin{tabular}{lr}
\toprule
Quantity & Value \\
\midrule
Epoch counts tested & 3 and 5 \\
Total simulation phases & 60 \\
Original three-epoch direct evaluations & 615 \\
Three-epoch refinement direct evaluations & 256 \\
Five-epoch direct evaluations & 1287 \\
Unique closeout direct evaluations & 2158 \\
Feasible policies found & 0 \\
Best $\phi$ & 5.83075 \\
Best active constraint & WG RAP \\
\bottomrule
\end{tabular}
\end{table}

The best observed schedule was a three-epoch refined policy:
\begin{center}
\begin{tabular}{lrrr}
\toprule
Control & Epoch 1 & Epoch 2 & Epoch 3 \\
\midrule
Annual intake & 147 & 144 & 122 \\
Retention rate & 0.5774 & 0.5696 & 0.5468 \\
UTE & 20.0 & 20.0 & 19.4119 \\
PAA & 30 & 30 & 30 \\
Maximum manning percentage & 170.324 & 176.677 & 182.110 \\
FLUG quota per phase & 2 & 2 & 3 \\
IPUG quota per phase & 0 & 0 & 0 \\
\bottomrule
\end{tabular}
\end{center}

The best policy satisfied the final total-pilot constraint, but it failed the
minimum total-pilot window after the assessment start. It also left WG RAP and
FL RAP violations. Table \ref{tab:best-violations} lists the positive
violations at this nearest miss.

\begin{table}[htbp]
\centering
\caption{Positive constraint relaxations required by the best observed policy.}
\label{tab:best-violations}
\begin{tabular}{lr}
\toprule
Constraint & Positive violation \\
\midrule
Total-pilot window & 533 pilots \\
WG RAP shortfall & 5.83075 \\
FL RAP shortfall & 2.94844 \\
\bottomrule
\end{tabular}
\end{table}

The active constraint distribution over unique closeout direct evaluations was
1128 FL RAP, 731 WG RAP, and 299 total-pilot window. The minimum observed value
of each core
constraint was negative somewhere in the run: $-2331$ for the total-pilot
window, $-0.9335$ for WG RAP, $-0.8856$ for FL RAP, and $-9.8083$ for IP RAP.
Thus, the search did find policies that satisfied each core constraint in
isolation. It did not find a policy that satisfied the inventory window, WG RAP,
and FL RAP constraints at the same time.

The observed relaxation study reached the same nearest miss. The minimum
L-infinity normalized relaxation over total-pilot window, WG RAP, FL RAP, and
IP RAP was 5.83075 at schedule \texttt{refine\_0008}. This is evidence from the
executed search, not a proof that no feasible policy exists.

\section{Local Control Authority Diagnostic}

A finite-difference diagnostic was run around the best schedule using
perturbations equal to 5 percent of each configured control range, with integer
controls rounded and perturbations clipped at bounds. This is a local
linearization diagnostic only; it is not a controller design or a substitute
for direct verification.

The strongest local sensitivities are summarized in Table
\ref{tab:authority}. Negative sensitivity means increasing the control reduces
the corresponding violation locally.

\begin{table}[htbp]
\centering
\caption{Dominant local control authority near the best trajectory.}
\label{tab:authority}
\begin{tabular}{llr}
\toprule
Response & Highest-authority control & Sensitivity \\
\midrule
$\phi$ & Epoch 3 IPUG quota & 1.920 \\
Total-pilot window & Epoch 1 retention rate & 90.91 \\
WG RAP & Epoch 2 retention rate & 1.587 \\
FL RAP & Epoch 2 retention rate & 6.005 \\
IP RAP & Epoch 2 retention rate & -3.679 \\
\bottomrule
\end{tabular}
\end{table}

The refined diagnostic is more cautionary than a simple ``turn one knob''
story. Increasing IPUG from zero worsens the current objective and FL RAP near
the best point. Retention has strong authority, but the sign depends on epoch
and response: it helps IP RAP in epoch 2 while worsening WG RAP, FL RAP, and the
total-pilot window locally. The diagnostic therefore points to a
joint-feasibility problem, not to a single constraint that is impossible to
satisfy on its own.

\section{Limitations}

A 2158-evaluation structured search does not prove global infeasibility. It
tests two policy classes: three-epoch and five-epoch open-loop schedules with
bounded controls and surrogate-guided candidate proposal. A fully free
per-phase schedule would have many more degrees of freedom. That option was
deliberately deferred to keep the implementation interpretable and the evidence
tractable.

The GPR surrogate is only a candidate generator. Its fitted length scales
suggest weak sensitivity in several sampled directions, but those diagnostics
are not feasibility evidence. Direct physics remains the source of truth.

\section{Recommendations}

The recommended next steps are:
\begin{enumerate}
  \item Keep the physics backend with worker-parallel direct evaluation as the
  baseline for final claims.
  \item Continue with structured open-loop searches before adding an MPC or a
  broader controls architecture.
  \item For additional search, move only with a justified new policy class or
  targeted requirement change; the tested five-epoch class did not improve the
  nearest miss.
  \item For decision support, start the requirements trade study from the best
  observed miss: 533 pilots in the total-pilot window, 5.83075 WG RAP, and
  2.94844 FL RAP.
\end{enumerate}

\section{Conclusion}

The viability workflow now has a physics-backed search and diagnostic loop for
structured dynamic policies. The implemented dynamic-policy model treats policy
design as finite-horizon control while staying close to the existing simulator
interfaces. The primary result is negative: 2158 unique direct physics
evaluations found no feasible policy in the tested three- and five-epoch
structured open-loop classes. The run still produced a clear nearest miss,
quantified the required relaxations, and identified local control authority
around the binding constraints. The next step is a requirements trade study or a
new policy class justified by these near-miss diagnostics.

\appendix
\section{Commands and Validation}

The primary dynamic search was run with the following command:
\begin{verbatim}
RUN=outputs/viability/dynamic_policy_search/run_3epoch_512_32768_096
venv/bin/python -m src.viability.cli dynamic-search \
  --config outputs/viability/dynamic_policy_search/config_physics.yaml \
  --output-dir ${RUN} \
  --epochs 3 \
  --initial-samples 512 \
  --optimizer-pool-size 32768 \
  --verify-top 96 \
  --workers 8 \
  --checkpoint-every 25
\end{verbatim}

The targeted three-epoch refinement was run with:
\begin{verbatim}
venv/bin/python -m src.viability.cli dynamic-refine \
  --config outputs/viability/dynamic_policy_search/config_physics.yaml \
  --previous-evaluations outputs/viability/dynamic_policy_search/run_3epoch_512_32768_096/all_evaluations.parquet \
  --diagnostic-sensitivity outputs/viability/dynamic_policy_search/run_3epoch_512_32768_096/diagnostic/local_sensitivity.csv \
  --output-dir outputs/viability/dynamic_policy_search/run_3epoch_refine_2048_131072_256 \
  --epochs 3 \
  --local-samples 2048 \
  --optimizer-pool-size 131072 \
  --verify-top 256 \
  --workers 8 \
  --checkpoint-every 25
\end{verbatim}

The five-epoch open-loop search was run with:
\begin{verbatim}
venv/bin/python -m src.viability.cli dynamic-search \
  --config outputs/viability/dynamic_policy_search/config_physics.yaml \
  --output-dir outputs/viability/dynamic_policy_search/run_5epoch_1024_131072_256 \
  --epochs 5 \
  --initial-samples 1024 \
  --optimizer-pool-size 131072 \
  --verify-top 256 \
  --workers 8 \
  --checkpoint-every 25
\end{verbatim}

The refined-best local diagnostic was run with:
\begin{verbatim}
venv/bin/python -m src.viability.cli dynamic-diagnose \
  --config outputs/viability/dynamic_policy_search/config_physics.yaml \
  --evaluations outputs/viability/dynamic_policy_search/run_3epoch_refine_2048_131072_256/all_evaluations.parquet \
  --output-dir outputs/viability/dynamic_policy_search/run_3epoch_refine_2048_131072_256/diagnostic \
  --epochs 3 \
  --perturbation-fraction 0.05 \
  --workers 8 \
  --checkpoint-every 10
\end{verbatim}

The relaxation study was run with:
\begin{verbatim}
venv/bin/python -m src.viability.cli dynamic-relaxation-study \
  --config outputs/viability/dynamic_policy_search/config_physics.yaml \
  --evaluations outputs/viability/dynamic_policy_search/run_3epoch_512_32768_096/all_evaluations.parquet \
  --evaluations outputs/viability/dynamic_policy_search/run_3epoch_refine_2048_131072_256/refinement_evaluations.parquet \
  --evaluations outputs/viability/dynamic_policy_search/run_5epoch_1024_131072_256/all_evaluations.parquet \
  --output-dir outputs/viability/dynamic_policy_search/relaxation_study_v1 \
  --top-n 50
\end{verbatim}

The branch validation reported:
\begin{itemize}
  \item 66 viability tests passed with \texttt{venv/bin/python -m unittest}.
  \item \texttt{venv/bin/python -m compileall src/viability viability\_dashboard.py} passed.
  \item \texttt{git diff --check} passed.
  \item Generated outputs under \path{outputs/viability} remained ignored and
  untracked.
\end{itemize}

\begin{thebibliography}{9}
\raggedright

\bibitem{internal-audit}
\textit{Internal Phase Backend Audit and Surrogate Plan}.
Repository document:
\path{docs/internal_phase_surrogate_audit_plan.md}.

\bibitem{physics-performance}
\textit{Viability Physics Evaluation Performance Plan}.
Repository document:
\path{docs/viability_physics_performance_plan.md}.

\bibitem{dynamic-status}
\textit{Dynamic Policy / Finite-Horizon Control Status}.
Ignored run artifact:
\path{outputs/viability/dynamic_policy_search/dynamic_control_status.md}.

\bibitem{dynamic-report}
\textit{Dynamic Policy / Finite-Horizon Control Search}.
Ignored diagnostic artifact:
\path{outputs/viability/dynamic_policy_search/run_3epoch_refine_2048_131072_256/diagnostic/dynamic_control_report.md}.

\end{thebibliography}

\end{document}
